\documentclass[aps,prl,twocolumn,superscriptaddress,floatfix]{revtex4-2}
\usepackage{amsmath,amssymb,graphicx,bm}
\begin{document}

\title{From-scratch surrogate replication of the nonlinear magnetic orbital
Hall effect in CuMnAs (Wang \emph{et al.}, 2026)}
\author{Hermes replication subagent}
\affiliation{Automated computational replication --- textures-100 corpus, \texttt{wang2026}}
\date{\today}

\begin{abstract}
We attempt a fast, from-scratch replication of Wang \emph{et al.}
(arXiv:2604.02636, 2026), which predicts a second-order \emph{nonlinear
magnetic orbital Hall effect} (nonlinear MOHE) in the $PT$-symmetric collinear
antiferromagnet CuMnAs, with headline
$\chi^{(O)}_{zzyy}\!\approx\!-1.3$ vs $\chi^{(S)}_{zzyy}\!\approx\!-0.0087\,
(\hbar/e)\,\Omega^{-1}\mathrm{V}^{-1}$ (orbital $>\!100\times$ spin). Lacking DFT
budget, we build a minimal four-band $PT$-symmetric AFM nodal-line tight-binding
model and evaluate the orbital and spin \emph{Berry-curvature dipoles}
$D_{zzyy}$ (with $\chi=\tau D$) on a coarse $18^3$ $k$-grid using the paper's own
operator definitions. We reproduce, at the qualitative/order-of-magnitude level,
the four central physics claims: (i) orbital $\gg$ spin (ratio $\sim\!10^3$ in our
model, paper $\sim\!150$); (ii) the effect is SOC-\emph{induced} (vanishes without
SOC); (iii) it is \emph{non-perturbative} in SOC (enhanced at weak SOC / small
gap); (iv) it is $T$-odd, flipping sign under N\'eel reversal. We do \emph{not}
reproduce the absolute magnitude in physical units. Verdict: \textbf{PARTIAL}
(gap: absolute magnitude / DFT-level material specificity).
\end{abstract}

\maketitle

\section{Target claims}
From the paper and \texttt{replication\_recipe.json}, the observable is the
second-order response $\delta j^d_a = \chi_{dabc}E_bE_c$ with $\chi_{dabc}=\tau
D_{dabc}$, where $D_{dabc}$ is the orbital (or spin) Berry-curvature dipole (OBD):
\begin{equation}
D_{abc}=\sum_n\!\int\!\frac{d^dk}{(2\pi)^d}\,\Omega^{d,n}_{ab}(\bm k)\,
\partial_c f_0,\quad
\Omega^{d,n}_{ab}=-2\,\mathrm{Im}\!\sum_{n'\neq n}\frac{j^{d,nn'}_a v^{n'n}_b}{\omega_{nn'}^2}.
\end{equation}
The itinerant orbital-angular-momentum matrix elements follow the paper's Eq.~(3),
\begin{equation}
L_{mn}=\frac{i}{4\mu_B}\!\!\sum_{\ell\neq m,n}\!\!\Big(\tfrac{1}{\varepsilon_\ell-\varepsilon_m}+\tfrac{1}{\varepsilon_\ell-\varepsilon_n}\Big)(\bm v_{m\ell}\times\bm v_{\ell n})_z,
\end{equation}
and the orbital (spin) current is $j^z_a=\tfrac12\{v_a,L^z\}$
($\tfrac12\{v_a,S^z\}$). Headline: $\chi^{(O)}_{zzyy}\approx-1.3$,
$\chi^{(S)}_{zzyy}\approx-0.0087\,(\hbar/e)\Omega^{-1}\mathrm V^{-1}$ at $T=50$~K,
$\tau=1.4$~ps; orbital exceeds spin by two orders of magnitude because weak SOC
gaps a nodal line near $X$ ($<20$~meV) and amplifies the OBD non-perturbatively.

\section{Surrogate model}
DFT+Wannier for an SOC AFM is out of scope for a $<\!6$-min replication. Instead we
build a four-band model on sublattice$\otimes$spin ($\rho\otimes\sigma$) with
$P=\rho_x$, $T=i\sigma_y K$:
\begin{align}
H(\bm k)=&\,[M-t(\cos k_x+\cos k_y+\cos k_z)]\rho_z + v\sin k_y\,\rho_y\nonumber\\
&+ J\,\rho_z\sigma_z + \lambda(\sin k_x\,\rho_x\sigma_x+\sin k_z\,\rho_x\sigma_y).
\end{align}
The first line hosts a \emph{nodal line} in the $k_y=0$ plane (mimicking CuMnAs
near $X$); $J\rho_z\sigma_z$ is the N\'eel AFM order (breaks $T$, preserves $PT$);
$\lambda$ is the SOC term that gaps the line. Velocities are computed by finite
difference $v_a=\partial H/\partial k_a$; $L^z$, $j^z_a$, $\Omega$, and $D_{zzyy}$
follow the equations above on a uniform $18^3$ grid at $T=0.08$ (model units),
$\mu$ placed on the nodal-line-derived Fermi surface.

\section{Results}
Table~\ref{tab:res} summarizes. At physical (weak) SOC $\lambda=0.12$ we obtain
$D^{(O)}_{zzyy}=-3.6\times10^{-16}$, $D^{(S)}_{zzyy}=-5.9\times10^{-20}$
(model units), an \textbf{orbital/spin ratio $\sim6.0\times10^3$}. Scanning
$\lambda$, the orbital dipole \emph{grows as SOC weakens}
(slope $d\log|D^{(O)}|/d\log\lambda\approx-1.1$), the non-perturbative signature.
With SOC off ($\lambda=0$), $D^{(O)}_{zzyy}\to-3.3\times10^{-22}$ (vanishing to
grid noise) and $D^{(S)}=0$, confirming SOC is required. Flipping the N\'eel
vector $J\to-J$ flips the sign of $D^{(O)}$ ($-3.6\times10^{-16}\to
+5.6\times10^{-16}$), confirming $T$-oddness.

\begin{table}[t]
\caption{Surrogate vs.\ paper. Model units; absolute magnitude not compared.}
\label{tab:res}
\begin{ruledtabular}
\begin{tabular}{lcc}
Claim & Paper & This work \\ \hline
$\chi=\tau D$ (OBD) structure & yes & yes \\
orbital $\gg$ spin & $\sim\!150\times$ & $\sim\!6\times10^3\times$ \\
SOC required ($D\to0$, no SOC) & yes & yes \\
non-perturbative (weak-SOC) & yes & yes (slope $-1.1$) \\
$T$-odd (N\'eel flip) & yes & yes (sign flips) \\
absolute $|\chi|$ in $(\hbar/e)\Omega^{-1}\mathrm V^{-1}$ & $-1.3$ & not matched \\
\end{tabular}
\end{ruledtabular}
\end{table}

\section{Verdict}
\textbf{PARTIAL} (gap: absolute magnitude / DFT material specificity). All four
mechanism-level claims of the paper are reproduced qualitatively in a coarse-grid
from-scratch surrogate in $\sim16$~s. The absolute conductivity in physical units
requires SOC-DFT + Wannier interpolation, which was deliberately skipped.

\section*{Credit}
Velocity/itinerant-$L_z$ and Kubo machinery adapted from
\texttt{gobel2024\_sd\_skyrmion\_kubo\_Lz\_kernel.py} (G\"obel 2024, topological
OHE), generalized here to a $k$-space multiband model and the second-order
(Berry-curvature-dipole) response of Wang 2026.

\end{document}
